\section{Lecture 8: Model Selection, $R^2$, and Multicollinearity}

\subsection{Review of Distributions}

\[
F = \frac{\chi^2_m / m}{\chi^2_n / n} \sim F_{m,n},
\]

\[
T = \frac{Z}{\sqrt{\chi^2_n / n}} \sim t_n, \quad Z \sim N(0,1),
\]

\[
\chi^2_n = \sum_{i=1}^n Z_i^2, \quad Z_i \sim N(0,1),
\]

\[
t_n^2 \sim F_{1,n}.
\]

\subsection{SLR F-Test Review}

\textbf{Model:} $Y = \beta_0 + \beta_1 X + \varepsilon$.

\textbf{t-test:} $H_0: \beta_1 = 0 \quad H_1: \beta_1 \neq 0$.
\[
T = \frac{\bhat{\beta}_1 - \beta_1}{\hat{\sigma}/\sqrt{s_{xx}}} \sim t_{n-2}.
\]

\textbf{F-test:} $H_0: \beta_1 = 0 \quad H_1: \beta_1 \neq 0$.
\[
F = \frac{\SSreg/1}{\RSS/(n-2)} \sim F_{1,n-2}.
\]

\textbf{Relationship:} $T_{\text{observed}}^2 = F_{\text{observed}}$.

\subsection{Partial F-Test}

Test whether a subset of $k$ predictors should be jointly added to a model that already includes the other $p - k$ predictors.

\textbf{Full model:}
\[
Y = \beta_0 + \beta_1 x_1 + \cdots + \beta_p x_p + \varepsilon.
\]

\textbf{Reduced model:}
\[
Y = \beta_0 + \beta_{k+1} x_{k+1} + \cdots + \beta_p x_p + \varepsilon.
\]

\textbf{Hypothesis:}
\[
H_0: \beta_1 = \beta_2 = \cdots = \beta_k = 0,
\]
where $k < $ number of predictors.
\[
H_1: \text{at least one } \beta_j \neq 0.
\]

For $p$ predictors, there are $2^p$ models that can be constructed.

\begin{keyresult}[Comparison of Models]
\[
\SST_{\text{full}} = \SST_{\text{reduced}},
\]
\[
\RSS_{\text{reduced}} \ge \RSS_{\text{full}},
\]
\[
\mathrm{SS}_{\mathrm{reg,reduced}} \le \mathrm{SS}_{\mathrm{reg,full}}.
\]

If $\RSS_{\text{reduced}} \gg \RSS_{\text{full}}$, prefer the full model.
If $\RSS_{\text{reduced}} \approx \RSS_{\text{full}}$, prefer the reduced model.
\end{keyresult}

\subsubsection{Degrees of Freedom}

\[
df_{\text{full}} = n - p - 1,
\]

\[
df_{\text{reduced}} = n - (p - k) - 1,
\]

\[
df_{\text{difference}} = k.
\]

\subsubsection{Test Statistic}

\begin{keyresult}[Partial F-Test Statistic]
\[
F = \frac{(\RSS_{\text{reduced}} - \RSS_{\text{full}})/k}{\RSS_{\text{full}}/(n - p - 1)} \sim F_{k, n-p-1}.
\]
\end{keyresult}

\textbf{Decision rule:}
\begin{itemize}
\item The test rejects $H_0$ at level $\alpha$ if $\text{p-value} < \alpha$ or if
\[
F_{\text{obs}} \geq F_{\alpha,\,k,\,n-p-1}.
\]
\begin{itemize}
\item There exists a statistically significant linear relationship between $Y$ and at least one of the $k$ predictors.
\item We prefer the full model.
\end{itemize}

\item The test fails to reject $H_0$ at level $\alpha$ if $\text{p-value} > \alpha$ or if
\[
F_{\text{obs}} < F_{\alpha,\,k,\,n-p-1}.
\]
\begin{itemize}
\item There does not exist a statistically significant linear relationship between $Y$ and any of the $k$ predictors.
\item We prefer the reduced model.
\end{itemize}
\end{itemize}

The number of possible models for $p$ predictors is
\[
\text{Number of models} = 2^p.
\]

\subsection{Case Study}

Assume a response variable $Y$ and two predictors $X_1, X_2$:
\[
Y = \beta_0 + \beta_1 X_1 + \beta_2 X_2 + \varepsilon.
\]

\textbf{Step 1: t-Tests}. We perform individual t-tests:
\[
H_0: \beta_1 = 0, \qquad H_0: \beta_2 = 0.
\]

Suppose the p-value for $\beta_1$ is large (p-value $= 0.88$).

\textbf{Step 2: F-Test}. Test overall model significance:
\[
H_0: \beta_1 = \beta_2 = 0
\quad \text{vs} \quad
H_1: \beta_1 \neq 0 \text{ or } \beta_2 \neq 0.
\]

Assume the p-value is very small ($p\text{-value} = 2.5 \times 10^{-6}$). Thus, the null model is not preferred.

\textbf{Step 3: Partial F-Test}. Compare nested models:
\[
\text{Model 1 (reduced): } Y = \beta_0 + \beta_2 X_2 + \varepsilon \quad (\beta_1 = 0),
\]
\[
\text{Model 2 (full): } Y = \beta_0 + \beta_1 X_1 + \beta_2 X_2 + \varepsilon.
\]

\begin{itemize}
\item If p-value is large, prefer the reduced model.
\item If p-value is small, prefer the full model.
\end{itemize}

\subsection{Coefficient of Determination}

\begin{defn}{$R^2$}{def:r2}
$R^2$ measures the proportion of variation in the response variable $Y$ that is explained by the model:
\[
R^2 = \frac{\SSreg}{\SST} = 1 - \frac{\RSS}{\SST}.
\]
\end{defn}

\[
0 \le R^2 \le 1.
\]

\subsubsection{Interpretation of $R^2$ (Multiple Regression)}

If $R^2 = 1$, then $\RSS = 0$. All observed points lie exactly on the fitted regression hyperplane, and the covariates $X_1, X_2, \dots, X_p$ account for all variation in $Y$.

When $\bhat{\beta}_1 = \bhat{\beta}_2 = \cdots = \bhat{\beta}_p = 0$, then $R^2 = 0$. The model explains no variation in the response variable; equivalently, there is no linear association between $X_1, X_2, \dots, X_p$ and $Y$ in the sample data.

The closer $R^2$ is to 1, the stronger the linear association between the covariates $X_1, X_2, \dots, X_p$ and the response variable $Y$.

\subsection{Coefficient of Correlation}

The sample correlation coefficient ($r$) for $n$ pairs $(x_1, y_1), \dots, (x_n, y_n)$ is defined as
\[
r = \frac{\sum_{i=1}^{n} (x_i - \bar{x})(y_i - \bar{y})}
{\sqrt{\sum_{i=1}^{n} (x_i - \bar{x})^2}\sqrt{\sum_{i=1}^{n} (y_i - \bar{y})^2}}
= \frac{S_{xy}}{\sqrt{S_{xx}}\sqrt{S_{yy}}}.
\]

\[
-1 \leq r \leq 1.
\]

In simple linear regression,
\[
r = \operatorname{sign}(\bhat{\beta}_1)\sqrt{R^2}.
\]

\subsection{Coefficient of Correlation: Proof that $r^2 = R^2$}

The sign is clear; it remains to verify that $r^2 = R^2$.

\[
r^2 = \frac{\left(\sum_{i=1}^{n} (x_i - \bar{x})(y_i - \bar{y})\right)^2}
{\sum_{i=1}^{n} (x_i - \bar{x})^2 \sum_{i=1}^{n} (y_i - \bar{y})^2}
= \frac{\bhat{\beta}_1 \sum_{i=1}^{n} (x_i - \bar{x})(y_i - \bar{y})}{\SST}.
\]

\[
\bhat{\beta}_1 \sum_{i=1}^{n} (x_i - \bar{x})(y_i - \bar{y})
= \bhat{\beta}_1 \left(\mathbf{x} - \frac{1}{n}\bX\mathbf{x}\right)^\top
\left(\mathbf{y} - \frac{1}{n}\bX\mathbf{y}\right)
\]

\[
= \bhat{\beta}_1 \mathbf{x}^\top \left(\mathbf{I} - \frac{1}{n}\bX\right)\mathbf{y}
\]

\[
= (\bhat{\beta}_0 \mathbf{1} + \bhat{\beta}_1 \mathbf{x} - \bhat{\beta}_0 \mathbf{1})^\top
\left(\mathbf{I} - \frac{1}{n}\bX\right)\mathbf{y}
\]

\[
= (\mathbf{H}\mathbf{y})^\top \left(\mathbf{I} - \frac{1}{n}\bX\right)\mathbf{y}
- \bhat{\beta}_0 \mathbf{1}^\top \left(\mathbf{I} - \frac{1}{n}\bX\right)\mathbf{y}
\]

\[
= \mathbf{y}^\top \left(\mathbf{H} - \frac{1}{n}\bX\right)\mathbf{y}
= \SSreg.
\]

\textbf{Result:}
\[
r^2 = \frac{\bhat{\beta}_1 \sum_{i=1}^{n} (x_i - \bar{x})(y_i - \bar{y})}{\SST}
= \frac{\SSreg}{\SST} = R^2.
\]

\begin{figure}[H]
    \centering
    \includegraphics[width=0.6\linewidth]{pictures/1.png}
\end{figure}

\subsection{Cautions When Interpreting $R^2$}

\begin{caution}
There are no general rules about how large $R^2$ should be.

It is possible to have a large value of $R^2$ even when the linear regression model is not ideal.

$R^2$ is not meaningful if a linear regression model is not appropriate.

It is possible to obtain very high $R^2$ by fitting overly complicated models in multiple linear regression.
\end{caution}

\subsubsection{Model Comparison Using $R^2$}

\textbf{Models:}
\[
\text{Model 1: } Y = \beta_0 + \beta_1 X_1 + \varepsilon,
\]
\[
\text{Model 2: } Y = \beta_0 + \beta_1 X_1 + \beta_2 X_2 + \varepsilon,
\]
\[
\text{Model 3: } Y = \beta_0 + \beta_1 X_1 + \beta_2 X_2 + \beta_3 X_3 + \varepsilon.
\]

\textbf{Always:}
\[
R^2(\text{Model 3}) \geq R^2(\text{Model 2}) \geq R^2(\text{Model 1}).
\]

\begin{caution}
Thus, $R^2$ cannot be used to compare models with different numbers of predictors. Use adjusted $R^2$ instead.
\end{caution}

\subsection{Adjusted $R^2$}

\textbf{Recall:}
\[
R^2 = \frac{\SSreg}{\SST} = 1 - \frac{\RSS}{\SST},
\qquad R^2 \in [0,1].
\]

As the number of covariates increases, $R^2$ also increases. Therefore, $R^2$ is not suitable for comparing different models. To compare models with different numbers of predictors, use $R^2_{\text{adj}}$ instead of $R^2$.

\begin{keyresult}[Adjusted $R^2$]
The adjusted $R^2$ is defined as
\[
R^2_{\text{adj}} = 1 - \frac{\RSS/(n - p - 1)}{\SST/(n - 1)}
= 1 - \frac{n - 1}{n - p - 1}\frac{\RSS}{\SST}.
\]

\[
-\frac{p}{n - p - 1} \leq R^2_{\text{adj}} \leq R^2.
\]
\end{keyresult}

\subsection{Example: Computing $R^2$ and $R^2_{\text{adj}}$}

\begin{figure}[H]
    \centering
    \includegraphics[width=0.5\linewidth]{pictures/2.png}
\end{figure}

\textbf{Given:}
\[
\SSreg = 352.87, \quad \RSS = 111.46, \quad n = 150, \quad p = 1.
\]

\[
R^2 = \frac{\SSreg}{\SST}
= \frac{352.87}{352.87 + 111.46}
\approx 0.76.
\]

\[
R^2_{\text{adj}} = 1 - \frac{n - 1}{n - p - 1}\cdot \frac{\RSS}{\SST}
= 1 - \frac{150 - 1}{150 - 1 - 1} \cdot \frac{111.46}{352.87 + 111.46}
\approx 0.68.
\]

\textbf{Interpretation:} Approximately 76\% of the variation in $Y$ is explained by the model. After adjusting for the number of predictors, about 68\% of the variation remains explained. This indicates a reasonably strong fit, though the drop from $R^2$ to $R^2_{\text{adj}}$ reflects the penalty for model complexity.

\subsection{Multicollinearity}

Consider a multiple linear model with $p$ predictors:
\[
Y = \beta_0 + \beta_1 x_1 + \cdots + \beta_p x_p + \varepsilon,
\qquad \varepsilon \sim N(0,\sigma^2).
\]

The least squares estimator is
\[
\bbeta = (\bX^\top \bX)^{-1} \bX^\top \bY.
\]

Under $\varepsilon \sim N_n(0,\sigma^2 I)$,
\[
\bbeta \sim N_{p+1}\!\left(\bbeta,\, \sigma^2 (\bX^\top \bX)^{-1}\right).
\]

For $j = 0,1,\dots,p$,
\[
\bhat{\beta}_j \sim N\!\left(\beta_j,\, \sigma^2 (\bX^\top \bX)^{-1}_{j+1,j+1}\right),
\qquad
\Var(\bhat{\beta}_j) = \sigma^2 (\bX^\top \bX)^{-1}_{j+1,j+1}.
\]

\begin{caution}
\boxed{\Var(\bhat{\beta}_j) \text{ becomes large when other predictors are highly linearly correlated with } x_j.}
\end{caution}

\subsection{Two Predictors and Variance Inflation}

Consider
\[
Y = \beta_0 + \beta_1 x_1 + \beta_2 x_2 + \varepsilon.
\]

Let $r_{12}$ be the correlation between $x_1$ and $x_2$, and let $S_{x_j}$ be the standard deviation of $x_j$. Then for $j=1,2$,
\[
\Var(\bhat{\beta}_j)
= \frac{1}{1 - r_{12}^2}\cdot \frac{\sigma^2}{(n-1) S_{x_j}^2}.
\]

When $|r_{12}|$ is large, $\Var(\bhat{\beta}_j)$ becomes large.

\subsubsection{Conclusion}

Multicollinearity increases the variance of coefficient estimates, making them unstable and less reliable.

\subsection{Variance Inflation Factor (VIF)}

Correlation among predictors increases the variance of the estimated regression coefficients by a factor of
\[
\frac{1}{1 - r_{12}^2}.
\]

\subsubsection{Example (Two Predictors)}

When $r_{12} = 0.99$ (highly correlated),
\[
\Var(\bhat{\beta}_j)
= \frac{1}{1 - 0.99^2}\cdot \frac{\sigma^2}{(n-1)S_{x_j}^2}
= 50.25 \cdot \frac{\sigma^2}{(n-1)S_{x_j}^2}.
\]

When $r_{12} = 0$ (no correlation),
\[
\Var(\bhat{\beta}_j)
= \frac{1}{1 - 0^2}\cdot \frac{\sigma^2}{(n-1)S_{x_j}^2}
= \frac{\sigma^2}{(n-1)S_{x_j}^2}.
\]

\subsection{General Case: $p$ Predictors}

Consider
\[
Y = \beta_0 + \beta_1 x_1 + \beta_2 x_2 + \cdots + \beta_p x_p + \varepsilon.
\]

To measure how well $x_j$ is explained by the other predictors, define
\[
x_j = \alpha_0 + \sum_{k \ne j} \alpha_k x_k + \varepsilon.
\]

Let $R_j^2$ be the coefficient of determination from this regression. Then
\[
\Var(\bhat{\beta}_j)
= \frac{1}{1 - R_j^2}\cdot \frac{\sigma^2}{(n-1)S_{x_j}^2},
\qquad j = 1,2,\dots,p.
\]

\subsection{Variance Inflation Factor}

\begin{defn}{Variance Inflation Factor}{def:vif}
\[
\mathrm{VIF}(\bhat{\beta}_j) = \frac{1}{1 - R_j^2}.
\]

It measures how much the variance of $\bhat{\beta}_j$ is inflated due to multicollinearity.
\end{defn}

\subsection{Rule of Thumb}

\begin{caution}
\[\boxed{
\mathrm{VIF}(\bhat{\beta}_j) > 5 \;\;\Rightarrow\;\; \text{possible severe multicollinearity}
\quad (\text{often } >10 \text{ is considered strong}).}
\]
\end{caution}

\begin{notebox}
\textbf{Note:} $R_j^2$ measures how well $X_j$ can be predicted by the other predictors. $R_j^2 \neq$ overall $R^2$.
\end{notebox}

\clearpage
